\documentclass[11pt,a4paper]{article}
\usepackage[utf8]{inputenc}
\usepackage[T1]{fontenc}
\usepackage[spanish,es-tabla]{babel}
\usepackage{geometry}
\usepackage{graphicx}
\usepackage{float}
\usepackage{hyperref}
\usepackage{tabularx}
\usepackage{booktabs}
\usepackage{array}
\geometry{margin=2.5cm}
\hypersetup{colorlinks=true,linkcolor=blue,urlcolor=blue}
\setlength{\parindent}{0pt}
\setlength{\parskip}{0.55em}
\newcolumntype{Y}{>{\raggedright\arraybackslash}X}
\newcommand{\code}[1]{\texttt{\detokenize{#1}}}
\newcommand{\evidencia}[2]{%
\begin{figure}[H]\centering
\IfFileExists{#1}{\includegraphics[width=0.88\linewidth]{#1}}{%
\fbox{\begin{minipage}[c][3cm][c]{0.82\linewidth}\centering Pendiente de captura: \texttt{\detokenize{#1}}\end{minipage}}}
\caption{#2}
\end{figure}}

\begin{document}
\begin{titlepage}
\centering
\vspace*{1cm}
{\Large Pontificia Universidad Javeriana}\\
{\large Fundamentos de Ingeniería de Software}\\[1.5cm]
{\Huge \textbf{DriveControl / AutoMinder Enterprise}}\\[0.4cm]
{\Large Equipo SYNTIX TECH}\\[1.2cm]
{\LARGE \textbf{Milestone 4 - Dashboard, alertas y cierre del sistema}}\\[1.2cm]
\begin{tabular}{ll}
Profesor: & Luis Gabriel Moreno Sandoval\\
Fecha: & 19 de mayo de 2026\\
Repositorio: & \url{https://github.com/puj-course/FIS_2610_3517_G4}\\
\end{tabular}
\vfill
\end{titlepage}

\tableofcontents
\newpage

\section{Resumen ejecutivo del milestone}
El Milestone 4 cerró dashboard, alertas, CI/CD, Docker, SonarCloud, SMS/Twilio, métricas y evidencia final. GitHub live registra 132 issues cerradas, 20 HU cerradas estrictas y 3 PRs asociados, con el PR \#628 abierto.

\section{Objetivos del milestone}
\begin{itemize}
\item Consolidar dashboard y alertas documentales.
\item Evidenciar Docker, CI/CD, SonarCloud y SMS.
\item Centralizar la entrega final de rúbrica.
\item Cerrar postmortem técnico y de proceso.
\end{itemize}

\section{Alcance funcional}
El alcance incluye cierre del sistema, métricas de calidad, despliegue, SMS, evidencia y documentación. No incluye releases/tags porque no existen verificables.

\section{Evidencia ágil}
\begin{tabularx}{\linewidth}{l r Y}
\toprule
Elemento & Valor & Evidencia\\
\midrule
HU cerradas & 20 & Issues con patrón estricto de HU.\\
Issues cerradas & 132 & Milestone cerrado para issues.\\
PRs & 3 & PRs \#232, \#488 y \#628.\\
Commits & 277 & Ventana S11-S13 en Git local.\\
\bottomrule
\end{tabularx}

\evidencia{../../Entrega-Final/evidencias/05_github_issues_sprint13.png}{Issues finales de Sprint 13.}
\evidencia{../../Entrega-Final/evidencias/06_github_prs_finales.png}{PRs finales.}
\evidencia{../../Entrega-Final/evidencias/01_agile_milestones_semestre.png}{Milestones del semestre.}

\section{Métricas del milestone}
\begin{tabularx}{\linewidth}{Y r}
\toprule
Métrica & Valor\\
\midrule
HU cerradas & 20\\
Issues cerradas & 132\\
Issues abiertas & 0\\
Commits por ventana de sprint & 277\\
PRs asociados & 3\\
PRs mergeadas & 2\\
PRs abiertas & 1\\
\bottomrule
\end{tabularx}

Participación por issues asignadas: \code{samuelfl680} 56, \code{Sarm-m} 43, \code{solonlosada2006} 30, \code{juanvargax} 22 y \code{Juserora} 17.

\section{Evaluación del producto}
Se produjo el cierre funcional y documental del sistema con evidencia de calidad, despliegue, SMS y metodología. El PR \#628 queda explícitamente pendiente.

\section{Evaluación del proceso}
El proceso se enfocó en cierre de rúbrica, PRs, CI/CD y evidencia. Funcionó la respuesta asíncrona ante bloqueos, pero faltó planear releases/tags.

\section{Evaluación del esfuerzo}
El esfuerzo fue alto: 277 commits en S11-S13 y 132 issues cerradas. La participación estuvo concentrada en cierre técnico, DevOps, QA y documentación.

\section{Postmortem del milestone}
\subsection{Qué salió bien}
Se centralizó evidencia final, se documentó CI/CD, Docker, SonarCloud y SMS, y se cerraron issues críticas.

\subsection{Qué salió mal}
Quedó abierto el PR \#628, no hay releases/tags y parte de la evidencia se recolectó bajo presión.

\subsection{Causas raíz}
La causa raíz principal fue tratar algunas validaciones funcionales y evidencias como actividades de cierre, no como parte del DoD permanente.

\subsection{Impacto}
El impacto es una integración pendiente y trazabilidad por versión limitada.

\subsection{Acciones correctivas}
Resolver o justificar PR \#628, documentar ausencia de releases y exigir prueba funcional local antes de PR.

\subsection{Acciones preventivas y responsables sugeridos}
QA Lead debe custodiar prueba funcional; Configuration Manager debe definir versión; Product Owner debe documentar datos demo; DevOps debe vigilar SonarCloud temprano.

\section{Retrospectiva estrella de mar}
\subsection{Comenzar a hacer}
\begin{itemize}
\item Definir política de releases o justificar su ausencia.
\item Ejecutar pruebas funcionales locales antes de PRs de cierre.
\item Documentar datos demo coherentes con el usuario activo.
\end{itemize}

\subsection{Más de}
\begin{itemize}
\item Monitorear SonarCloud durante todo el sprint.
\item Revisar PRs con criterio funcional.
\item Capturar evidencia al validar cada flujo.
\end{itemize}

\subsection{Seguir haciendo}
\begin{itemize}
\item Mantener trazabilidad entre issues, PRs, commits y capturas.
\item Usar Docker Compose para reproducibilidad.
\item Responder colaborativamente ante bloqueos de cierre.
\end{itemize}

\subsection{Menos de}
\begin{itemize}
\item Acumular evidencia para los últimos días.
\item Depender del pipeline como única validación.
\item Usar datos demo sin reglas de consistencia.
\end{itemize}

\subsection{Dejar de hacer}
\begin{itemize}
\item Abrir PRs de cierre sin validación funcional.
\item Ignorar hotspots o coverage hasta el final.
\item Presentar releases si no son verificables.
\end{itemize}

\section{Plan de mejora final}
\begin{tabularx}{\linewidth}{Y Y Y Y}
\toprule
Acción & Responsable & Evidencia esperada & Criterio de éxito\\
\midrule
Resolver PR pendiente & QA Lead / equipo & PR cerrado o justificado & Cierre sin ambigüedad\\
Definir política de versión & Configuration Manager & Tag/release o decisión & Trazabilidad clara\\
Validar antes de PR & QA Lead & Checklist & Menos regresiones\\
Documentar datos demo & Product Owner & Guía de datos & Alertas reproducibles\\
\bottomrule
\end{tabularx}

\section{Conclusión del milestone}
M4 cerró el sistema con evidencia técnica y metodológica suficiente. Los pendientes se declaran sin inventar datos: PR \#628 abierto y ausencia de releases/tags verificables.

\section{Anexos o evidencias pendientes}
\begin{itemize}
\item Pendiente: cierre o justificación de PR \#628.
\item No existen releases/tags verificables.
\item Pendiente de exportación desde GitHub API: reviews completas por PR.
\end{itemize}
\end{document}
